
\subsection{Social-Level Theories}

Although interest in Marxist-Leninist theories has waned, there has been a tremendous growth of research on the relationships between regime type, particularly democratic regimes, and war (analyzed by Ray in this volume). Scholars have also devoted attention to the "diversionary" use of force for domestic political purposes, to the impact of ethnonationalism on international conflict, and to other domestic sources of international conflict.

"rally around the flag" effect (Mueller 1973).

This is the age-old "scapegoat hypothesis" or "diversionary theory of war" (Levy 1989a).

As Bodin argued, "the best way of preserving a state, and guaranteeing it against sedition, rebellion, and civil war is to...find an enemy against whom [the subjects] can make common cause" (quoted in Levy 1989a, p. 259).

research on the scapegoat hypothesis

The absence of significant findings regarding the relationship between internal and external conflict contrasts sharply with evidence of external scapegoating from historical and journalistic accounts and from a growing body of case-study evidence.

Most of these studies have focused on democratic regimes, partly because of the ease of measuring elite support levels but also because of the assumption that the greater political accountability of leaders in democratic states makes them more likely than authoritarian leaders to engage in external scapegoating.

These models help to explain "gambling for resurrection" through risky foreign policies by poor leaders who expect to be removed by their constituents (Downs \& Rocke 1995).

This raises the question of how external actors perceive and respond to diversionary actions.

Marxist-Leninist models of imperialism, like models of scapegoating, implicitly assume that external expansion and the use of force serve the interests of the ruling class or elite but not those of society as a whole.

If the costs of aggressive foreign policies are too high---and we have enough examples of "imperial overstretch" to suggest that they sometimes are---it is not clear how elites maintain their power. As Snyder (1991) argued, groups concentrated enough to benefit from overexpansion are too narrow to control state policy, while those broad enough to control policy are too diffuse to reap benefits from overexpansionist policies.

For Snyder, this rationalist model of logrolling is not sufficient to explain overexpansion or the stability of the ruling coalition.

\subsection{Individual-Level Theories}
decision-making is characterized by “bounded rationality”
\textbf{Assumptions:}
\begin{enumerate}
    \item [(a)] that external and internal structures and social forces are not translated directly into foreign policy choices;
    \item [(b)] that key decision-makers vary in their definitions of state interests, assessments of threats to those interests, and/or beliefs as to the optimum strategies to achieve those interests; and
    \item [(c)] that differences in the content of actors’ belief systems, in the psychological processes through which they acquire information and make judgments and decisions, and in their personalities and emotional states are important intervening variables in explaining observed variation in state behaviors with respect to issues of war and peace.
\end{enumerate}

\textbf{Prospect Theory} one of the most recent attempts to apply a social-psychological model to international relations (Kahneman \& Tversky (1979)).

Prospect theory assumes that people evaluate outcomes in terms of deviations from a reference point or aspiration level rather than in terms of net-asset levels.

This asymmetry of behavior with respect to gains and losses means that the way people frame their reference point in any given choice problem is critical.

\subsection{General Trends}

\begin{itemize}
    \item shift away from the systemic level: dissatisfaction with the failure of structural systemic models to explain enough of the variance in war and peace. This shift has led to rising interest in both dyadic-level behavior and societal-level explanatory variables.
    \item Whatever the relationship between concentrations of military and economic capabilities at the systemic level, there is substantial evidence that at the dyadic level an equality of capabilities is significantly more likely to lead to war.
    \item In the past decade most international relations theorists have moved beyond their earlier preoccupation with explanations at a single level of analysis and debates about which level provides the most powerful explanations. They now give much more attention to interaction effects among variables at different levels in the processes leading to war
\end{itemize}
